\scanpage{80}
% archive foliation: 79
on veut faire abstraction des questions de représentabilité sauf pour les $\Gamma^{\nu_i}(X_i)$)
\[\Gamma^{\nu_1}(X_1)\otimes\cdots\otimes\Gamma^{\nu_r}(X_r)\longleftarrow\Gamma^\nu(X_1,\ldots,X_r)\]
on aura un hom \unsure{fonctoriel en} $M$ dans $\unsure{\mathcal P|P}_1$
\[\unsure{\mathcal P|P}_\nu(X_1,\ldots,X_r;M)\longleftarrow \unsure{\mathcal P|P}_{1,\ldots,1}(\Gamma^{\nu_1}X_1,\ldots,\Gamma^{\nu_r}X_r;M)\]
par composition avec les applications canoniques de $\unsure{\mathcal P|P}$
\diagram{\begin{tikzcd}
X_1\arrow[r,"\unsure{e_{\nu_1}|\gamma_1}"]&\Gamma^{\nu_1}X_1&\cdots&X_r\arrow[r,"\unsure{e_{\nu_r}|\gamma_r}"]&\Gamma^{\nu_r}(X_r)
\end{tikzcd}}
et on suppose que c'est un \inserted{iso}. Moyennant ceci, on a
\[\Gamma^n(X_1\times\cdots\times X_r)\simeq\sum_{\substack{\nu=(\nu_i)\in\mathbb N^r\\\sum\nu_i=n}}\underbrace{\Gamma^{\nu_1}(X_1)\otimes\cdots\otimes\Gamma^{\nu_r}(X_r)}_{=\Gamma^{\nu_1,\ldots,\nu_r}(X_1,\ldots,X_r)}\]
[NB Il suffirait de \unsure{libeller} ces hypothèses relatives $r=\unsure{n|2}$ --- pour $r=2$ d'abord, et les déduire ensuite pour $r$ quelc.

En particulier on trouve, pour $n=1$, en considérant le produit $X\times0$, pour $1=\Gamma^0(0)$
\[\deleted{\Gamma^1(X\times1)\simeq}\]
\[X\otimes1\simeq\Gamma^1(X)\otimes\Gamma^0(0)\simeq\underbrace{\Gamma^{1,0}(X,0)\simeq X}_{\text{car }\unsure{\mathcal P|P}_{1,0}(X,0;M)\simeq \unsure{\mathcal P|P}_1(X,M)}\]
et donc $1\otimes X\simeq X$. \deleted{\unsure{N.B. on a aussi}} \deleted{$\unsure{X\otimes1}$}

On suppose qu'il y a des produits tensoriels stricts dans les \unsure{$\unsure{\mathcal P|P}_1$ en jeux}, et que les $\Gamma^nX$ satisfont à \unsure{la condition} plus haut.
\margin{\unsure{formant une} \unsure{catégorie} \unsure{rigide}}

